\begin{abstract}
Efficient 3D segmentation decoders are commonly evaluated by aggregate accuracy
and computational cost, leaving a basic question unanswered: \emph{where does
additional decoder computation actually improve a prediction?}
We study this question through a controlled reimplementation of a full 3D
UX-Net decoder and its compressed EffiDec3D counterpart on 13-organ abdominal
CT segmentation.  Rather than treating every disagreement as a benefit, we
define voxel-wise decoder marginal utility using positive transitions
(efficient wrong, full correct), negative transitions (efficient correct, full
wrong), and their net difference.  Our current study yields three consistent
observations.  First, errors are spatially heterogeneous: boundary error is
$3.93\times$ interior error, and organ difficulty is negatively associated with
organ size ($\rho=-0.544$).  Second, the aggregate benefit of the full decoder
is small and bidirectional; across two efficient-decoder seeds, positive and
negative transition rates differ by only $0.065$--$0.130\%$ of evaluated
voxels.  Third, this marginal benefit is predictable: subject-level
entropy--net-utility correlations are 0.665 and 0.646 across the two seeds,
while a preliminary voxel oracle places approximately 86\% of positive
transitions in the top 20\% entropy-ranked foreground voxels.
These findings motivate selective decoder allocation, but do not by themselves
demonstrate realizable savings.  We therefore introduce a paired,
subject-bootstrap evaluation protocol based on net utility and contiguous 3D
regions, separating oracle opportunity from deployable selection.
\todo{Replace the preliminary voxel-oracle sentence with corrected region-level
O9 results and add cross-backbone/cross-dataset findings before submission.}
\end{abstract}
